\subsection{Historical Origins of Operations Research}\label{sec:introduction:historical-origins-of-operations-research}

This section is short and mostly historical, but it explains \textbf{why OR emerged exactly when it did}.

\highspace
\begin{flushleft}
    \textcolor{Green3}{\faIcon{fighter-jet} \textbf{OR during World War II}}
\end{flushleft}
Operations Research developed as a recognizable discipline during \textbf{World War II}. Military organizations faced complex operational problems involving scarce resources, uncertainty, and many possible decisions. Teams of scientists were asked to study how operations could be conducted more efficiently.

\highspace
Typical problems included:
\begin{itemize}
    \item Where to place radar systems;
    \item How to organize convoys;
    \item How to allocate limited military resources;
    \item How to improve logistics.
\end{itemize}
The basic idea was already the same as modern OR:
\begin{equation*}
    \text{Real operational problem}
    \rightarrow
    \text{quantitative analysis}
    \rightarrow
    \text{better decision}
\end{equation*}
The word \textbf{operations} in \emph{Operations Research} comes from this origin. \hl{Researchers were literally studying military operations.} The important conceptual point is that OR was born from a need to make \textbf{better use of scarce resources under complex conditions}.

\highspace
\begin{flushleft}
    \textcolor{Green3}{\faIcon{industry} \textbf{Expansion to industry and management}}
\end{flushleft}
After the war, the techniques developed for military applications became available for civilian use. At the same time, companies and organizations were becoming much larger and more complex.

\highspace
This created problems such as:
\begin{itemize}
    \item Production planning;
    \item Transportation;
    \item Inventory management;
    \item Workforce scheduling;
    \item Logistics;
    \item Supply-chain coordination.
\end{itemize}
So methods originally designed for military operations naturally transferred to industrial and managerial problems. The same underlying question remained: \emph{\textbf{how should limited resources be allocated to achieve the best result?}} For example:
\begin{equation*}
    \text{military vehicles}
    \longrightarrow
    \text{delivery vehicles}
\end{equation*}
\begin{equation*}
    \text{troop deployment}
    \longrightarrow
    \text{workforce scheduling}
\end{equation*}
\begin{equation*}
    \text{military logistics}
    \longrightarrow
    \text{supply-chain logistics}
\end{equation*}
The applications changed, but the mathematical structure was often similar. This is why OR became closely associated with \textbf{Management Science}. The two terms are often used almost interchangeably, although \emph{Management Science} tends to emphasize applications to managerial decision-making.

\highspace
\begin{flushleft}
    \textcolor{Green3}{\faIcon{microchip} \textbf{Role of algorithms and computing}}
\end{flushleft}
OR could expand rapidly because two technological developments happened at the same time. First, there was strong progress in \textbf{optimization methods and numerical algorithms}. Second, \textbf{computers became increasingly available}. These two developments reinforce each other. A mathematical optimization model may have an enormous number of possible solutions.

\highspace
For example, with $n$ binary decisions:
\begin{equation*}
    x_i\in\{0,1\}
\end{equation*}
There are potentially $2^n$ different combinations. For $n=100$:
\begin{equation*}
    2^{100}
    \approx 1.27\times 10^{30}
\end{equation*}
So just having a mathematical formulation is not enough. We need:
\begin{equation*}
    \text{Efficient algorithms}
    +
    \text{Computing power}
\end{equation*}
Without efficient algorithms, computers would simply enumerate possibilities too slowly. Without computers, even good algorithms would be limited to very small problems. So modern OR emerged from the combination:
\begin{equation*}
    \text{Mathematical modeling}
    +
    \text{Algorithmic advances}
    +
    \text{Computers}
\end{equation*}
